\begin{solapartat}
\punts{0,25} $|\vec{E}| = \dfrac{\Delta V}{d} = \dfrac{15000}{0,015} = 10^{6}\ \mathrm{V/m}$

\punts{0,5} Per moure els electrons cap a la peça a soldar, la força ha d'anar dirigida del càtode a l'ànode, i com que els electrons tenen càrrega negativa el camp elèctric té el sentit oposat de la força:

\figura[0.55]{sol_fig1.png}

\punts{0,5} El potencial alt correspon a l'ànode.

Justificació: el camp elèctric sempre apunta en la direcció que disminueix el potencial.

\destaca{Alternativament}, atès que les partícules han d'accelerar-se sota l'acció del camp, llavors el treball fet pel camp elèctric ha de ser positiu:
\[ W_\mathrm{Camp} = -q\Delta V = |e|\Delta V = |e|(V_\text{ànode} - V_\text{càtode}) > 0 \;\Rightarrow\; V_\text{ànode} > V_\text{càtode} \]
\end{solapartat}

\begin{solapartat}
\punts{0,4} $\Delta E_C = \dfrac{1}{2}mv^2 - 0 = W_\mathrm{Camp} = -q\Delta V = |e|\Delta V = 2,403\times 10^{-15}\ \mathrm{J}$

\punts{0,4} $v = \sqrt{\dfrac{2W_\mathrm{Camp}}{m}} = 7,26\times 10^{7}\ \mathrm{m/s}$

\punts{0,45} $v = 7,26\times 10^{7}\ \mathrm{m/s}$. Atès que la diferència de potencial és zero entre la peça i l'ànode, el camp elèctric no fa cap treball o també es pot afirmar que el camp elèctric en aquesta regió és zero i, per tant, l'electró es mou a velocitat constant.
\end{solapartat}
